% §3 Threat Model & Problem Definition — target 0.5 page
% Purpose: formalize. ACSAC reviewers expect this explicitly.

\section{Threat Model and Problem Definition}
\label{sec:threat}

\paragraph{Network-scale pentest}
We define \emph{network-scale penetration testing} as the task of: (i)~discovering
all live hosts in a target IP range, (ii)~enumerating services and identifying
vulnerabilities on each host, (iii)~producing evidence of exploitability where
feasible, and (iv)~aggregating findings at the network level. The \emph{unit of
evaluation} is the network, not an individual host.

\paragraph{Attacker model}
The attacker holds network-layer access to the target subnet (post-ingress,
pre-compromise). They hold no prior credentials, no privileged vantage, and no
knowledge of deployed software beyond what is observable through scans. The
attacker has access to standard offensive tooling (Nmap, MQTT clients, SSH, \ldots)
and an LLM with tool-calling. Time and budget are bounded.

\paragraph{Goals}
\begin{enumerate}
  \item Enumerate all devices reachable from the entry point.
  \item For each device, identify vulnerabilities from misconfiguration, weak
  credentials, information disclosure, insecure protocols, and known CVEs.
  \item For each finding, attempt to produce \emph{evidence} at one of three
  levels: \textsf{L1}~detection (port open, banner leak), \textsf{L2}~exploitation
  (authenticated session, protocol handshake), \textsf{L3}~data exfiltration
  (configuration files, credentials, sensor data).
  \item Report findings at the network level (severity-ranked, deduplicated).
\end{enumerate}

\paragraph{Scope and non-goals}
We evaluate at the \emph{IP/application layer} on deployed IoT networks. We
do \emph{not} evaluate: (a)~physical/side-channel attacks on hardware,
(b)~RF-layer attacks requiring SDR captures,
(c)~firmware extraction from physical devices,
(d)~post-exploitation persistence or impact,
(e)~insider threats or supply-chain compromise.
Hardware-assisted attacks are left as future work (\cref{sec:conclusion}).
